\begin{frame}{Full Training Workflow}
\begin{itemize}
    \item \textbf{Step 1: Initial Setup}
    \begin{itemize}
        \item Start with a pretrained model and just finetune when possible; it saves time and improves results.
        \item Define an initial architecture without regularization (e.g., dropout) or augmentations.
        \item Set up validation strategy and choose an appropriate evaluation loss and metric.
        \item Train the model to get a \textbf{baseline score}.
    \end{itemize}
\end{itemize}
\end{frame}

\begin{frame}{Full Training Workflow}
\begin{itemize}
    \item \textbf{Step 2: Improvement Process}
    \begin{itemize}
        \item Overfitting is common at the start; use regularization techniques like:
        \begin{itemize}
            \item Dropout, batch normalization,...
        \end{itemize}
        \item Perform error analysis to identify weaknesses and choose appropriate augmentations or preprocessing.
        \item Tune hyperparameters: layers, epochs, learning rate, batch size, etc.
        \item Optionally, use ensembling to boost scores (requires more resources).
    \end{itemize}
    \item \textbf{Key Tip:} Track scores and improvements at every step to measure progress effectively.
\end{itemize}
\end{frame}

\begin{frame}{Measuring an Improvement Against the Baseline}
\begin{itemize}
    \item The baseline from Step 1 is the thing every later run is compared to. Freeze it and keep its checkpoint.
    \item Change \textbf{one} thing at a time, and hold the seed fixed across both arms, so the difference you measure is the change and not the draw.
    \item Before believing a gain, find out how big the noise is: rerun the baseline with a different seed and nothing else altered.
    \item If the improvement is smaller than that seed-only spread, it is not an improvement yet. Rerun both arms on three to five seeds and compare the means.
\end{itemize}
\end{frame}

\begin{frame}{Keep a Run Log}
\begin{itemize}
    \item Every run should leave a row behind: what you changed, what it scored, and enough information to reproduce it.
\end{itemize}
\vspace{0.5em}
\begin{center}
\small
\begin{tabular}{|c|c|c|l|c|}
    \hline
    \textbf{Run} & \textbf{Commit} & \textbf{Seed} & \textbf{Change from baseline} & \textbf{Val acc} \\
    \hline
    001 & a3f19c & 42 & baseline, no augmentation & 0.812 \\
    002 & a3f19c & 43 & baseline, seed only & 0.807 \\
    003 & 7d2e40 & 42 & + random crop and flip & 0.841 \\
    004 & 91b6aa & 42 & + dropout 0.3 in the head & 0.815 \\
    \hline
\end{tabular}
\end{center}
\vspace{0.5em}
\begin{itemize}
    \item Run 002 is the one that makes the table worth keeping: reseeding alone moved the score by $0.005$, so run 004 is noise while run 003 is real.
    \item A spreadsheet is enough for a course project, as long as you fill it in as you go rather than from memory afterwards.
    \item Tools such as TensorBoard, MLflow and Weights \& Biases automate this and version the artifacts too; the Introduction to AI course covers them under MLOps.
\end{itemize}
\end{frame}

\begin{frame}{Full Training Workflow}
\begin{itemize}
    \item \textbf{Step 3: Finalization}
    \begin{itemize}
        \item Save the optimized model for deployment.
        \item Use the model for inference in real-world applications.
    \end{itemize}
\end{itemize}
\end{frame} 